% Generated draft: requires mathematical and source review.
\documentclass[11pt]{article}
\usepackage[margin=25mm]{geometry}
\usepackage{fontspec}
\setmainfont{Latin Modern Roman}
\usepackage{amsmath,amssymb,mathtools,unicode-math}
\setmathfont{Latin Modern Math}
\usepackage{enumitem,needspace,xurl,hyperref}
\setlength{\parindent}{0pt}
\setlength{\parskip}{5pt}
\setlength{\emergencystretch}{3em}
\newcommand{\N}{\mathbb{N}}
\newcommand{\Z}{\mathbb{Z}}
\newcommand{\Q}{\mathbb{Q}}
\newcommand{\R}{\mathbb{R}}
\newcommand{\C}{\mathbb{C}}
\newcommand{\F}{\mathbb{F}}
\newcommand{\A}{\mathbb{A}}
\newcommand{\PP}{\mathbb{P}}
\newcommand{\E}{\mathbb{E}}
\newcommand{\eps}{\varepsilon}
\newcommand{\dd}{\,\mathrm{d}}
\begin{document}
\section*{UnsolvedMath ID: 30004701}
\textbf{Source problem code:} catalogue entry
\subsection{Definitions and mathematical statement}
A cyclic quadrilateral is a quadrilateral whose four vertices lie on a common circle in the Euclidean plane. The supplied record does not specify what quantities must be rational in the phrase ``rational cyclic quadrilaterals,'' nor does it specify the content of the reconstructions of Brahmagupta's theory or Euclid's lost porisms.

The supplied historical question is: \emph{How did Michel Chasles envision a common framework for reconstructing Brahmagupta's theory of rational cyclic quadrilaterals and Euclid's lost porisms?}

The requested object is an account of Chasles's intended framework. No definition of that framework, relevant primary-source passage, mathematical hypotheses, or criterion for a successful reconstruction is supplied. Consequently, this record does not determine a precise mathematical proposition with specified domains and quantifiers. The question is retained as a historical research question; its presupposition that Chasles envisioned such a common framework requires source verification.
\subsection{Short English statement}
Describe the common framework that the supplied question attributes to Michel Chasles for reconstructing Brahmagupta's rational cyclic quadrilaterals and Euclid's lost porisms. The intended framework and the meaning of ``rational'' are unspecified, so a precise mathematical statement cannot be recovered from this record alone.
\subsection{Sources}
\begin{itemize}
\item Supplied local UnsolvedMath registry record, permanent ID \(30004701\), source code \texttt{catalogue entry}, titled ``Chasles's Reconstruction of Brahmagupta and Euclid,'' in the Oberwolfach Reports Open Problems collection. The registry records proposed year \(2021\) and supplies no named proposer. Chasles is the subject of the historical inquiry; the record does not identify him as its proposer.
\item Supplied report citation: \emph{Mathematics and its Ancient Classics Worldwide: Translations, Appropriations, Reconstructions, Roles} (2022), Oberwolfach Reports, pp.~1347--1406, DOI \texttt{10.4171/owr/2021/26}. Supplied source and literature-source link, labeled ``Oberwolfach report'': \url{https://doi.org/10.4171/owr/2021/26}. The background describes Oberwolfach Reports as documenting workshops and their open-problem discussions. No question-specific page or primary Chasles text is supplied.
\item Literature assessment dated exactly \texttt{2026-08-22}: the supplied curation labels the question ``open'' and describes it as an underspecified historical question about a scholar's intended common framework, unsettled by the available workshop record. This is supplied dated curation, not an independently verified literature assessment; its time-sensitive status and cited claims require independent verification.
\item The supplied registry labels the statement \texttt{corrected\_verified} and reports a manual semantic, typographic, and LaTeX review asserting that the clean statement is self-contained and preserves its target. This is a supplied review claim; it does not resolve the formulation and source limitations described here.
\end{itemize}
\end{document}
